\section*{Capítulo \ref{ch:g11:scal} --- El producto escalar en el plano}

\begin{solution}{exo:g11:scal:1}
$(2,5)\cdot(3,-1) = 6 - 5 = 1$.

$(1,-4)\cdot(8,2) = 8 - 8 = 0$ (\omterm{def:g11:scal:orthogonal}{vectores ortogonales}).

$\vec u\cdot\vec v = 3 \times 4 \times \cos\frac\pi3 = 12 \times \frac12
= 6$.
\end{solution}

\begin{solution}{exo:g11:scal:2}
$\vect{AB}\cdot\vect{AC} = 6 \times 6 \times \cos 60^\circ = 18$.

Los \omterm{def:g11:vect:vector}{vectores} $\vect{AB}$ y $\vect{BC}$ forman un ángulo de
$180^\circ - 60^\circ = 120^\circ$ (colócalos con el mismo origen en
$B$), así que $\vect{AB}\cdot\vect{BC} = 36 \cos 120^\circ = -18$.
\end{solution}

\begin{solution}{exo:g11:scal:3}
$\vec u \cdot \vec v = 3(t - 8) + t = 4t - 24 = 0$ da $t = 6$.

$\vec u - \vec v = (t - 4,\ 0)$, luego
$\vec u \cdot (\vec u - \vec v) = t(t-4) + 0 = t^2 - 4t$, que se anula
para $t = 0$ o $t = 4$.
\end{solution}

\begin{solution}{exo:g11:scal:4}
$\vec u \cdot \vec v = 3 + 2 = 5$, $\norm{\vec u} = \sqrt5$,
$\norm{\vec v} = \sqrt{10}$, luego
\[
\cos\theta = \frac{5}{\sqrt5\,\sqrt{10}} = \frac{5}{5\sqrt2}
= \frac{\sqrt2}{2},
\]
así que $\theta = 45^\circ$ exactamente.
\end{solution}

\begin{solution}{exo:g11:scal:5}
$a^2 = 16 + 36 - 2 \times 4 \times 6 \times \cos\frac{2\pi}{3}
= 52 - 48 \times \left(-\frac12\right) = 76$, luego
$a = \sqrt{76} = 2\sqrt{19} \approx 8.7$.

Para el segundo triángulo, el teorema del \omterm{def:g11:trigo:cossin}{coseno} despejado en el ángulo
da
\[
\cos\widehat A = \frac{b^2 + c^2 - a^2}{2bc}
= \frac{25 + 16 - 49}{2 \times 5 \times 4} = -\frac{8}{40} = -\frac15,
\]
negativo: $\widehat A$ es obtuso.
\end{solution}

\begin{solution}{exo:g11:scal:6}
$\vect{AB}\,(4, -2)$ y $\vect{AC}\,(2, 4)$:
$\vect{AB}\cdot\vect{AC} = 8 - 8 = 0$, así que el ángulo en $A$ es recto.
Lados: $AB = \sqrt{16 + 4} = \sqrt{20}$,
$AC = \sqrt{4 + 16} = \sqrt{20}$ y, con $\vect{BC}\,(-2, 6)$,
$BC = \sqrt{40}$. Coherencia: $BC^2 = 40 = 20 + 20 = AB^2 + AC^2$, el
teorema del \omterm{def:g11:trigo:cossin}{coseno} con $\cos\widehat A = 0$ (Pitágoras).
\end{solution}

\begin{solution}{exo:g11:scal:7}
Forma centro--radio: $(x + 2)^2 + (y - 3)^2 = 25$.

Para la segunda circunferencia, completamos cuadrados:
\[
x^2 - 6x + y^2 + 2y = 6
\iff (x - 3)^2 - 9 + (y + 1)^2 - 1 = 6
\iff (x-3)^2 + (y+1)^2 = 16 :
\]
centro $(3, -1)$, radio $4$.
\end{solution}

\begin{solution}{exo:g11:scal:8}
$M(x,y)$ está en la circunferencia cuando $\vect{MA}\cdot\vect{MB} = 0$
con $\vect{MA}\,(-3 - x,\ -y)$ y $\vect{MB}\,(-x,\ 2 - y)$:
\[
(-3 - x)(-x) + (-y)(2 - y) = x^2 + 3x + y^2 - 2y = 0 .
\]
Para el origen: $0 + 0 + 0 - 0 = 0$, así que $O$ está en la
circunferencia; geométricamente, $\vect{OA}\,(-3,0)$ y $\vect{OB}\,(0,2)$
son \omterm{def:g11:scal:orthogonal}{ortogonales}, luego $O$ ve $[AB]$ bajo un ángulo recto. (Completando
cuadrados: centro $\left(-\frac32, 1\right)$, radio
$\frac{\sqrt{13}}{2}$.)
\end{solution}

\begin{solution}{exo:g11:scal:9}
La perpendicular por $P$ tiene como \omterm{def:g11:vect:vector}{vector} normal el \omterm{def:g11:vect:direction}{vector director} de
$d$, es decir, $(-1, 2)$: \omterm{def:g10:algebra:equation}{ecuación} $-x + 2y + c = 0$; pasando por
$P(1,3)$: $-1 + 6 + c = 0$, $c = -5$, luego $x - 2y + 5 = 0$.

Corte con $d$: de $2x + y = 7$, $y = 7 - 2x$; sustituyendo,
$x - 2(7 - 2x) + 5 = 5x - 9 = 0$, luego $x = \frac95$ e
$y = \frac{17}5$: $H\left(\frac95, \frac{17}5\right)$. Entonces
$\vect{PH}\left(\frac45, \frac25\right)$ y
\[
PH = \sqrt{\frac{16}{25} + \frac{4}{25}} = \frac{\sqrt{20}}{5}
= \frac{2\sqrt5}{5} \approx 0.89 .
\]
\end{solution}

\begin{solution}{exo:g11:scal:10}
Sumando
\[
\norm{\vec u + \vec v}^2 = \norm{\vec u}^2 + 2\vec u\cdot\vec v +
\norm{\vec v}^2
\quad\text{y}\quad
\norm{\vec u - \vec v}^2 = \norm{\vec u}^2 - 2\vec u\cdot\vec v +
\norm{\vec v}^2,
\]
los términos cruzados se cancelan y
$\norm{\vec u + \vec v}^2 + \norm{\vec u - \vec v}^2 = 2\norm{\vec u}^2 +
2\norm{\vec v}^2$. En un paralelogramo de lados $\vec u$ y $\vec v$, las
diagonales son $\vec u + \vec v$ y $\vec u - \vec v$: la suma de los
cuadrados de las dos diagonales es igual a la suma de los cuadrados de los
cuatro lados.
\end{solution}

\begin{solution}{exo:g11:scal:11}
\emph{1.} Como $I$ es el \omterm{prop:g10:coordgeom:midpoint}{punto medio}, $\vect{IB} = -\vect{IA}$ y
$IA = 2$. Entonces
\[
\vect{MA}\cdot\vect{MB}
= (\vect{MI} + \vect{IA})\cdot(\vect{MI} - \vect{IA})
= \norm{\vect{MI}}^2 - \norm{\vect{IA}}^2
= MI^2 - 4 .
\]

\emph{2.} La condición $\vect{MA}\cdot\vect{MB} = 12$ se convierte en
$MI^2 = 16$, es decir, $MI = 4$: el conjunto es la circunferencia de
centro $I$ y radio $4$.
\end{solution}

\begin{solution}{pb:g11:scal:1}
\textbf{1.} $\vec u \cdot \vec v = -3 + 8 = 5$; $\norm{\vec u} = 5$,
$\norm{\vec v} = \sqrt5$; $\cos\theta = \frac{5}{5\sqrt5} =
\frac{1}{\sqrt5}$: $\theta \approx 63.4^\circ$.

\textbf{2.} $6 - 4k = 0$: $k = \frac32$.

\textbf{3.} $W = 50 \times 10 \times \cos 60^\circ = 250$ julios: solo
trabaja la componente en la dirección del movimiento.

\textbf{4.} $c^2 = 25 + 49 - 2 \times 35 \times \frac12 = 39$:
$c = \sqrt{39}$.

\textbf{5.} $\cos\gamma = \frac{16 + 36 - 64}{2 \times 4 \times 6} =
-\frac14$: negativo, el ángulo es obtuso:
$\gamma \approx 104.5^\circ$.

\textbf{6.} Con \omterm{def:g10:coordgeom:system}{coordenadas}:
$\vec u \cdot \vec v = \cos a \cos b + \sin a \sin b$. Con geometría: son
dos \omterm{def:g11:vect:vector}{vectores} unitarios, luego
$\vec u \cdot \vec v = 1 \times 1 \times \cos(a - b)$. Igualando los dos
cálculos del mismo número:
$\cos(a - b) = \cos a \cos b + \sin a \sin b$.

\textbf{7.} $\cos(a + b) = \cos(a - (-b)) = \cos a \cos(-b) +
\sin a \sin(-b) = \cos a \cos b - \sin a \sin b$.

\textbf{8.} $\sin(a + b) = \cos\left(\frac\pi2 - a - b\right)
= \cos\left(\left(\frac\pi2 - a\right) - b\right)
= \cos\left(\frac\pi2 - a\right)\cos b +
\sin\left(\frac\pi2 - a\right)\sin b
= \sin a \cos b + \cos a \sin b$.

\textbf{9.} $\cos 15^\circ = \cos 45^\circ \cos 30^\circ +
\sin 45^\circ \sin 30^\circ = \frac{\sqrt2}{2} \cdot
\frac{\sqrt3}{2} + \frac{\sqrt2}{2} \cdot \frac12 =
\frac{\sqrt6 + \sqrt2}{4} \approx 0.966$. Y
$\sin 75^\circ = \cos 15^\circ$: el mismo valor.

\textbf{10.} Con $b = a$ en las preguntas 7 y 8:
$\cos 2a = \cos^2 a - \sin^2 a = 2\cos^2 a - 1 = 1 - 2\sin^2 a$ (usando
$\cos^2 + \sin^2 = 1$), y $\sin 2a = 2 \sin a \cos a$.

\textbf{11.} $\sin^2 15^\circ = \frac{1 - \cos 30^\circ}{2}
= \frac{2 - \sqrt3}{4}$, luego
$\sin 15^\circ = \frac{\sqrt{2 - \sqrt3}}{2}$. Elevando al cuadrado
$\frac{\sqrt6 - \sqrt2}{4}$:
$\frac{6 - 2\sqrt{12} + 2}{16} = \frac{8 - 4\sqrt3}{16} =
\frac{2 - \sqrt3}{4}$: el mismo número con dos disfraces.

\textbf{12.} $\vect{AP}$ se descompone en una parte paralela a $d$
(invisible para $\vec n$) y otra perpendicular a $d$, cuya longitud es la
distancia buscada; multiplicar escalarmente por la normal unitaria
$\frac{\vec n}{\norm{\vec n}}$ mata la primera parte y mide la segunda.
Con $A$ en $d$: $\vect{AP} \cdot \vec n = a x_0 + b y_0 -
(a x_A + b y_A) = a x_0 + b y_0 + c$, de donde la fórmula. Ejemplo:
$\frac{\abs{15 + 24 - 12}}{5} = \frac{27}{5} = 5.4$.

\textbf{13.} Radio $= 5.4$: $(x - 5)^2 + (y - 6)^2 = 29.16$.

\textbf{14.} Distancia $= \frac{\abs{-12}}{5} = 2.4$. Pie: la
perpendicular por el origen es $(3t, 4t)$; sobre la recta:
$9t + 16t = 12$, $t = \frac{12}{25}$: el pie es
$\left(\frac{36}{25}, \frac{48}{25}\right)$, a distancia
$5t = \frac{60}{25} = 2.4$. La fórmula y el cálculo honrado coinciden.

\textbf{15.} Base $AB = 6$ y altura $=$ distancia de $C$ a $(AB)$ (el eje
$y = 0$): $5$; área $15$. Recta $(AC)$: $5x - 2y = 0$, y
$\operatorname{dist}(B, (AC)) = \frac{\abs{30}}{\sqrt{29}}$; entonces
$\frac12 \times AC \times \frac{30}{\sqrt{29}} =
\frac12 \times \sqrt{29} \times \frac{30}{\sqrt{29}} = 15$: la misma área,
con cualquiera de las dos alturas.

\textbf{16.} $\operatorname{dist} = \frac{\abs{2 + 1 - 6}}{\sqrt2} =
\frac{3}{\sqrt2} \approx 2.12 > 2$: la recta pasa de largo sin tocar la
circunferencia; son exteriores.

\textbf{17.} Intercalamos $I$:
$\vect{MA} \cdot \vect{MB} = (\vect{MI} + \vect{IA}) \cdot
(\vect{MI} + \vect{IB}) = MI^2 + \vect{MI} \cdot (\vect{IA} +
\vect{IB}) + \vect{IA} \cdot \vect{IB}$. El término central muere
($\vect{IA} + \vect{IB} = \vec 0$) y
$\vect{IA} \cdot \vect{IB} = -IA^2 = -\frac{AB^2}{4}$ (\omterm{def:g11:vect:vector}{vectores} opuestos
de longitud $\frac{AB}2$).

\textbf{18.} $\vect{MA} \cdot \vect{MB} = 0 \iff MI^2 = \frac{AB^2}{4}
\iff MI = \frac{AB}{2}$: la circunferencia de centro $I$ y radio
$\frac{AB}2$, o sea, la de diámetro $[AB]$.

\textbf{19.} Desarrollando cada pareja desde $H$ ($\vect{BC} =
\vect{HC} - \vect{HB}$, etc.):
\[
\vect{HA} \cdot \vect{HC} - \vect{HA} \cdot \vect{HB}
+ \vect{HB} \cdot \vect{HA} - \vect{HB} \cdot \vect{HC}
+ \vect{HC} \cdot \vect{HB} - \vect{HC} \cdot \vect{HA} = 0 :
\]
todo se cancela por parejas; la identidad vale para \emph{cuatro} puntos
cualesquiera. Para nuestro $H$, los dos primeros \omterm{def:g11:scal:dot}{productos escalares} se
anulan por hipótesis, luego $\vect{HC} \cdot \vect{AB} = 0$: $H$ está
sobre la altura desde $C$. Tres alturas, un punto.

\textbf{20.} Las fórmulas de adición salieron de calcular un mismo
\omterm{def:g11:scal:dot}{producto escalar} con la cara de las \emph{\omterm{def:g10:coordgeom:system}{coordenadas}} y con la de las
\emph{\omterm{def:g11:scal:dot}{normas} y el ángulo}; la fórmula de la distancia, de la cara de la
\emph{proyección}; el \omterm{pb:g11:scal:1}{ortocentro}, de la cara de la \emph{bilinealidad}
(Chasles y \omterm{def:g10:algebra:expand}{desarrollar}). Calcula dos veces la misma cantidad, iguala y
cae un teorema: ese es todo el oficio del \omterm{def:g11:scal:dot}{producto escalar}.
\end{solution}
